\section{Related Works}
\label{sec:related_works}

\subsection{Vision-Language-Action (VLA) Models}

\noindent\textbf{Foundation VLAs.}
Vision-Language-Action (VLA) models~\cite{brohan2022rt,brohan2023rt,team2024octo,bjorck2025gr00t,li2025hamster,black2024pi_0,yang2025magma,pertsch2025fast,driess2025knowledge,bu2025agibot,team2025gemini,wang2025vla} have recently emerged as a promising paradigm for embodied AI by training vision-language backbones on large-scale robot demonstrations. Works such as OpenVLA~\cite{kim24openvla} and $\pi_0$~\cite{black2024pi_0} achieve language-conditioned manipulation through end-to-end policy learning, while Magma~\cite{yang2025magma} co-trains on heterogeneous human and robot data. HAMSTER~\cite{li2025hamster} and TraceVLA~\cite{zheng2024tracevla} further leverage 2D visual trajectories to boost spatial-action connections. Despite success on routine manipulation, these imitation-based approaches struggle with long-horizon planning and generalization to novel scenarios due to limited training data coverage.

\vspace{2mm}

\noindent\textbf{Reasoning VLAs.}
To overcome these limitations, recent works~\cite{zawalski2024robotic,zhao2025cot,lee2025molmoact,wu2025you,qu2025eo,huang2025thinkact,kim2025robot,yuan2025embodied,abdolmaleki2025gemini} integrate explicit reasoning mechanisms into VLA architectures. Supervised approaches~\cite{zawalski2024robotic,zhao2025cot,lee2025molmoact,qu2025eo} introduce intermediate reasoning through chain-of-thought annotations. Embodied CoT~\cite{zawalski2024robotic} and Hi-Robot~\cite{shi2025hi} synthesize reasoning labels via pretrained foundation models. To perform vision-centric reasoning~\cite{man2025argus,sarch2025grounded} beyond pure text, CoT-VLA~\cite{zhao2025cot} employs visual goal generation and MolmoAct~\cite{lee2025molmoact} structures reasoning by spatial representations. Additionally, EO-1~\cite{qu2025eo} introduces interleaved vision-language-action pre-training to bridge reasoning and interaction. Recent works~\cite{yuan2025embodied,huang2025thinkact} alternatively leverage reinforcement fine-tuning to generate reasoning chains with designed rewards. Despite improved generalization, these reasoning VLAs suffer from high inference latency and inevitably introduce extraneous information that degrades action quality.


\subsection{Efficient Reasoning}

To address the inference latency of reasoning, recent LLM research explores various efficiency techniques~\cite{lee2025vlsi,dai2025stable,yuan2025efficient,xiang2025just,aggarwal2025l1,lee2025unified}. For example, RL-based approaches~\cite{dai2025stable,yuan2025efficient,xiang2025just,aggarwal2025l1} introduce length penalties to encourage shorter reasoning chains, though such methods can suffer from training instability. Beyond length control, latent reasoning methods~\cite{hao2024training,shen2025codi,zhang2025soft,cheng2024compressed,xu2025softcot} enable reasoning in continuous spaces, such as Coconut~\cite{hao2024training} using hidden states as continuous thoughts, CODI~\cite{shen2025codi} distilling explicit CoT into continuous space via teacher-student alignment, and Soft Thinking~\cite{zhang2025soft} generating weighted concept tokens. However, these LLM techniques cannot directly transfer to VLA tasks due to the need for spatial-temporal understanding and bridging semantic reasoning with embodied control.
Recently, ECoT-Lite~\cite{chen2025training} proposes reasoning dropout to accelerate embodied reasoning by skipping test-time reasoning traces. However, reasoning dropout can lead to inconsistent planning as it builds on supervised embodied CoT. Our proposed \ours{} distills reasoning into compact latent representations that naturally encode multimodal information, enabling robust reasoning-enhanced policy learning.